% =====================================================================
% 第4章：AR沉浸式可视化系统
% ChartForge 声明式图表生成框架
% =====================================================================
\chapter{AR沉浸式可视化系统}
\label{ch:ar_visualization}

为便于指称，本章将\textbf{图表意图形式化表示}简记为 \textbf{CIF}（Chart Intent Formalization），将\textbf{概率图表语法}简记为 \textbf{PCG}（Probabilistic Chart Grammar），将\textbf{语义-视觉对齐分数}简记为 \textbf{SVAS}（Semantic-Visual Alignment Score），将\textbf{多阶段生成精炼管线}简记为 \textbf{MS-GRP}（Multi-Stage Generative Refinement Pipeline），将\textbf{可组合图表代数}简记为 \textbf{CCA}（Composable Chart Algebra）。上述方法与第三章多模态交互、第五章实证分析共同构成完整的AR沉浸式分析系统。

\section{问题分析与研究动机}
\label{sec:ch4_motivation}

第三章实现了手势、语音与射线协同的多模态交互通道，为AR场景中的自然操作提供了输入基础。然而，第\ref{ch:ar_interaction}章解决的是”如何输入”的问题，尚未系统解决”呈现什么图表”以及”如何生成图表”的问题。现有AR可视化方案的研究已取得显著进展：Wang等\cite{wang2026ar_io_survey}从输入输出双视角系统梳理了AR人机交互技术全貌，Liu等\cite{Liu_Jia_Nie_Rosenberg_Interrante_Zhu-Tian_2025}综述了沉浸式分析中空间数据叙事与多视图联动的核心方法，Panzeri等\cite{panzeri2026cartographic}系统讨论了AR制图可视化的设计空间与评估框架，Chen等\cite{chen2025immersive_econ}探索了沉浸式环境在经济数据呈现中的应用潜力。尽管如此，上述方案在图表呈现方面仍存在三个根本性局限。其一，图表类型在设计阶段已固化编码——系统仅支持预设的省域热力图、时间轴折线和回归结果面板，用户无法根据分析需求动态生成新的图表类型或组合（如区域异质性分析时同时查看散点图与时序折线图），这要求开发人员提前为每种视图组合编写渲染逻辑。其二，从自然语言分析需求到图表库API调用序列之间存在显著的语义鸿沟——用户在AR穿戴场景下既无法便捷编程，也难以通过现有语音指令表达复杂的图表配置意图。其三，直接使用大语言模型生成图表代码面临典型的”生成即提交”问题——模型输出缺乏定量评估机制，而在AR可穿戴场景中，错误图表的误导成本远高于桌面场景。

\medskip
\noindent\textbf{基于React的声明式图表生成方法。}
React框架\cite{walke2013react}对前端开发的革命性贡献在于其核心洞察：与其手动管理DOM的增量变更，不如每次声明式地描述”UI应该长什么样”，由框架自动计算并应用最小变更集。这一范式的关键技术支柱包括Virtual DOM（内存中UI轻量表示）、Reconciliation（协调算法计算最小差异集）、声明式组件模型（$\text{UI} = f(\text{state})$）和单向数据流。

本文的核心洞察是：AR图表的生成可以类比UI渲染，将”AIGC生成AR图表”重新表述为”从意图描述到图表规约的声明式映射”。形式化地：

\begin{equation}
    \text{AR图表控件} = f(\text{用户意图},\; \text{数据模式},\; \text{设计约束})
    \label{eq:ch4_declarative}
\end{equation}

其中$f$是一个可学习的生成函数，由本章提出的ChartForge框架实现。该框架将图表生成从”手写渲染代码”的工程范式升级为”NL$\rightarrow$规约$\rightarrow$AR图表”的声明式AIGC范式，与第三章的多模态输入和第五章的计量数据资产协同工作。

本章的核心贡献包括：(C1)~提出图表意图形式化表示（CIF），将自然语言需求映射为结构化图表规约（\S\ref{sec:cif}）；(C2)~设计概率图表语法（PCG），以可学习的概率上下文无关文法描述图表生成空间（\S\ref{sec:pcg}）；(C3)~提出语义-视觉对齐分数（SVAS）作为生成质量的定量度量（\S\ref{sec:svas}）；(C4)~构建多阶段生成精炼管线（MS-GRP），通过”生成--评估--精炼”闭环迭代提升图表质量（\S\ref{sec:msgrp}）；(C5)~提出可组合图表代数（CCA），实现图表元素的声明式组合与复用（\S\ref{sec:cca}）。

% =====================================================================
\section{ChartForge系统架构}
\label{sec:ch4_architecture}

ChartForge系统架构采用四层组织，如图\ref{fig:ch4_architecture}所示。自底向上，意图层接收来自第三章多模态通道的自然语言查询，经CIF解析器输出结构化图表意图三元组$\langle\mathcal{D},\mathcal{V},\mathcal{I}\rangle$，与TSTQ--Fusion共享时间基准。生成精炼层以CIF为输入，经PCG束搜索生成候选图表规约，再经MS-GRP四阶段管线（粗生成$\rightarrow$语义验证$\rightarrow$视觉精炼$\rightarrow$交互注入）输出优化后的ChartSpec。渲染适配层将平台无关的ChartSpec映射为目标平台的图表渲染指令——AR端转换为Unity ChartRenderer可消费的JSON格式，Web端或离线评估则转换为Vega-Lite/ECharts格式以支持基线对比。交互运行时层管理用户与AR图表的交互事件，通过统一状态向量$\mathbf{s}_t$同步，确保ChartForge图表与省域地图、时间轴、结果面板在同一刷新周期内一致更新。

\begin{center}
  \includegraphics[width=0.72\textwidth]{ch4-1-architecture.png}
  \captionof{figure}{\label{fig:ch4_architecture}ChartForge系统架构：四层递进，NL输入经CIF解析与MS-GRP精炼后输出AR图表规约，与省域地图、时间轴、面板在统一状态向量驱动下协同呈现。}
\end{center}

ChartForge与React声明式UI的范式映射关系如表\ref{tab:react_map}所示。正如React的Virtual DOM作为平台无关的UI抽象，Chart Spec作为平台无关的图表抽象，使得同一份规约可渲染为Web（ECharts/Vega-Lite）、AR（Unity ChartRenderer）或桌面端（Plotly）等多种形态。

\begin{table}[H]
\centering
\caption{ChartForge与React的声明式范式映射}
\label{tab:react_map}
\begin{tabular}{l|l|l}
\toprule
\textbf{React 概念} & \textbf{ChartForge 对应} & \textbf{说明} \\
\midrule
Virtual DOM     & Chart Spec (中间规约)    & 图表的平台无关抽象表示 \\
Reconciliation  & MS-GRP 精炼管线         & 计算意图到图表的最优映射路径 \\
Component       & Chart Glyph / Layer     & 可复用的图表基本单元 \\
Props           & CIF 参数绑定             & 数据到视觉通道的属性传递 \\
Hooks           & Interaction Constraints & 声明式交互行为描述 \\
JSX             & Chart Definition DSL    & 图表声明语法 \\
Renderer        & Platform Adapter        & 跨平台（AR/Web/Desktop）图表渲染驱动 \\
\bottomrule
\end{tabular}
\end{table}

% =====================================================================
\section{图表意图形式化表示（CIF）}
\label{sec:cif}

CIF是ChartForge的中继表示层，借鉴React中Virtual DOM作为平台无关UI描述的角色，将模糊的自然语言需求$\mathcal{Q}$映射为结构化的图表意图三元组$\text{CIF}(\mathcal{Q}) = \langle \mathcal{D}, \mathcal{V}, \mathcal{I} \rangle$（式~\ref{eq:cif}）。三元组分别对应数据语义、视觉编码与交互约束三个正交维度，共同构成从”用户说了什么”到”系统需要画什么”的完整语义桥梁。

\begin{equation}\label{eq:cif}
    \text{CIF}(\mathcal{Q}) = \langle \mathcal{D},\; \mathcal{V},\; \mathcal{I} \rangle
\end{equation}

数据语义层$\mathcal{D} = \big\{ (f_i, t_i, s_i, a_i) \big\}_{i=1}^{n}$刻画用户意图中涉及的数据字段。对每个字段$f_i$（如\texttt{“DEL”}, \texttt{“province”}, \texttt{“year”}），CIF同时标注其数据类型$t_i$（离散、有序、数值或时序）、建议的视觉通道$s_i$（坐标轴、颜色、尺寸等）以及聚合函数$a_i$（求和、均值、计数或不聚合）。在省域经济分析场景中，数据语义层的字段词汇表与第五章指标体系严格对齐——省份、区域、年份、DEL、ES、PGDP、URBAN、INDS、TEIN等——确保了从用户说出的指标名到数据引擎中的列名再到回归表中的变量名的全程一致性。

视觉编码层$\mathcal{V} = \langle \mathcal{G}, \mathcal{E}, \Theta \rangle$指定图表的视觉呈现方式。$\mathcal{G} \in \Gamma$为目标图表类型，$\Gamma$覆盖柱状图、折线图、散点图、面积图、热力图、饼图、雷达图、桑基图、矩形树图、箱线图、仪表盘和漏斗图共12类常用图表；$\mathcal{E}: \mathcal{D} \to \{\text{position}, \text{color}, \text{size}, \text{shape}, \text{opacity}, \text{text}\}$将数据字段映射到视觉通道；$\Theta = \{\theta_1, \dots, \theta_k\}$收纳配色方案、字体、标注位置等样式参数。

交互约束层$\mathcal{I} = \big\{ (e_j, h_j, c_j) \big\}_{j=1}^{m}$描述用户对图表交互行为的期望，每条约束指定交互事件类型$e_j$（点击、悬停、刷选、缩放、筛选、下钻）、事件处理器规约$h_j$以及性能约束$c_j$（如响应时间低于200\,ms）。该层与第三章TSTQ--Fusion的交互事件模型共享相同的类型系统，确保交互事件可在ChartForge图表和AR视图之间无歧义路由。

CIF的解析由GPT-4o结合结构化提示词完成，系统提示词定义了完整的JSON输出模式，包含字段类型约束、图表类型枚举和交互事件白名单。解析后执行三项后处理：字段验证（检查解析出的字段名是否在领域词汇表中，对误识别字段进行模糊匹配修正）、类型推断补充（对缺失类型信息的字段从领域知识库中自动补充）以及置信度标记（解析失败的查询返回低置信度标记，由上游语音交互层通过AR HUD给出轻提示）。

% =====================================================================
\section{概率图表语法（PCG）}
\label{sec:pcg}

PCG是ChartForge的理论核心，借鉴了React Reconciliation中协调算法搜索最小变更集的思想——正如Reconciliation在Virtual DOM树上搜索最优更新路径，PCG在图表语法树上搜索最符合用户意图的图表派生路径。PCG形式上定义为一个五元组$\text{PCG} = \langle N, T, S, R, P \rangle$（式~\ref{eq:pcg}），其中各分量的含义与设计如下。

\begin{definition}[概率图表语法]
一个概率图表语法定义为一个五元组：
\begin{equation}\label{eq:pcg}
    \text{PCG} = \langle N,\, T,\, S,\, R,\, P \rangle
\end{equation}
\end{definition}

非终结符集合$N$表示图表的抽象组件层级，包含$\texttt{Chart}$（图表根节点）、$\texttt{Layout}$（布局容器）、$\texttt{Glyph}$（几何标记）、$\texttt{Axis}$（坐标轴）、$\texttt{Scale}$（比例尺）、$\texttt{Legend}$（图例）、$\texttt{Guide}$（辅助线）、$\texttt{Facet}$（分面）、$\texttt{Layer}$（图层）和$\texttt{Annotation}$（注释）共10个抽象组件。终结符集合$T$包含21种可渲染的图表原子元素，涵盖几何元素（$\texttt{BarGlyph}$、$\texttt{PointGlyph}$、$\texttt{LinePath}$、$\texttt{AreaPath}$、$\texttt{ArcGlyph}$、$\texttt{RectCell}$）、文字与刻度（$\texttt{TextLabel}$、$\texttt{TickMark}$、$\texttt{GridLine}$）、比例尺（$\texttt{ColorScale}$、$\texttt{SizeScale}$、$\texttt{OpacityScale}$、$\texttt{ShapeScale}$）、坐标轴（$\texttt{OrdinalAxis}$、$\texttt{QuantitativeAxis}$、$\texttt{TemporalAxis}$）以及辅助元素（$\texttt{LegendDef}$、$\texttt{GuideLine}$、$\texttt{FacetGrid}$、$\texttt{AnnotationText}$、$\texttt{AnnotationRegion}$）。起始符号$S = \texttt{Chart}$表示所有图表派生均从根节点开始。

产生式规则集$R$定义了从非终结符到符号串的合法展开方式。核心规则包括：图表根节点展开为布局、至少一个几何标记、可选辅助线与注释的组合（$\texttt{Chart} \to \texttt{Layout} \; \texttt{Glyph}^+ \; \texttt{Guide}^* \; \texttt{Annotation}^*$）；布局展开为横纵坐标轴与可选分面（$\texttt{Layout} \to \texttt{Axis}_x \; \texttt{Axis}_y \; \texttt{Facet}^*$）；几何标记可从六种原子图形中选择（$\texttt{Glyph} \to \texttt{BarGlyph} \mid \texttt{PointGlyph} \mid \ldots$）；坐标轴与比例尺各有相应的类型派生规则。当前PCG共定义23条产生式规则，构成一个紧凑但表达能力充分的图表生成空间——虽然规则数量有限，但通过递归组合可覆盖12类图表类型的任意合法实例。

概率分布$P: R \to [0, 1]$为每条产生式规则赋予可学习的权重，并满足对每个非终结符$A$其候选规则概率之和为1的归一化条件$\sum_{r \in R_A} P(r) = 1$。这一概率层使PCG区别于普通的上下文无关文法——在生成时不是随机选择规则，而是依据从图表语料库中学习到的先验偏好进行加权采样。

PCG的概率分布$P$通过图表语料库的最大似然估计学习得到。给定训练集$\mathcal{C} = \{(x_i, y_i)\}_{i=1}^{M}$（来自ChartIntent-10K数据集的7000条训练样本），训练目标为最大化图表语法树的似然$\mathcal{L}_{\text{PCG}} = \sum_{i=1}^{M} \sum_{r \in \text{parse}(y_i)} \log P(r \mid \text{parent}(r)) + \lambda \cdot \mathcal{R}(P)$，其中$\text{parse}(y_i)$是图表$y_i$对应的语法树展开所使用的产生式序列。为防止未见规则出现零概率，采用Laplace平滑为每条规则保留$1/(n+k)$的最小概率质量（$n$为规则使用总次数，$k$为该非终结符的候选规则数）。

\begin{equation}\label{eq:pcloss}
    \mathcal{L}_{\text{PCG}} = \sum_{i=1}^{M} \sum_{r \in \text{parse}(y_i)} \log P\big(r \mid \text{parent}(r)\big) + \lambda \cdot \mathcal{R}(P)
\end{equation}

在推理阶段，给定用户意图$\mathcal{Q}$的CIF表示，图表生成转化为在PCG语法树空间中求解最优派生路径的优化问题$T^* = \arg\max_{T \in \mathcal{T}(\text{PCG})} P(T \mid \text{PCG}) \cdot P(\text{CIF}(\mathcal{Q}) \mid T)$，其中$\mathcal{T}(\text{PCG})$是PCG可派生的所有语法树集合，两项因子分别编码了从语料库学习的先验图表偏好和当前用户意图的匹配度。本文采用束搜索（束宽$k=5$），在每个扩展步评估产生式规则的对数概率与CIF意图匹配分数之和。算法复杂度为$O(k \cdot |R| \cdot d)$，其中语法树最大深度$d \leq 10$，$|R| \approx 23$，每张图表生成仅需约1150次产生式扩展评估，在Python实现中单次采样耗时低于1\,ms。

\begin{equation}\label{eq:conditional}
    T^* = \arg\max_{T \in \mathcal{T}(\text{PCG})} \;
          \underbrace{P(T \mid \text{PCG})}_{\text{先验偏好}} \;\cdot\;
          \underbrace{P\big(\text{CIF}(\mathcal{Q}) \mid T\big)}_{\text{意图匹配度}}
\end{equation}

PCG具备两条基本的理论性质，为ChartForge的生成能力提供了形式化保障。
\begin{theorem}[完备性]\label{thm:completeness}
PCG的生成空间$\mathcal{T}(\text{PCG})$包含所有12类图表类型的任意合法实例。
\end{theorem}

\begin{proof}[证明思路]
对任意图表$C \in \Gamma$，其可分解为$(\text{Layout},\, \text{Glyph}^+,\, \text{Guide}^*,\, \text{Annotation}^*)$四元组。通过构造性证明：对每种图表类型$\gamma \in \Gamma$，存在至少一组产生式$R_\gamma \subseteq R$使得$S \Rightarrow^* C_\gamma$，且其概率$P(C_\gamma) > 0$。实验验证中，对12种图表类型各生成100个随机实例，所有类型均成功派生（覆盖率100\%）。
\end{proof}

\begin{theorem}[一致性]\label{thm:consistency}
PCG生成任意图表$C$的概率$P(C)$满足Chapman-Kolmogorov一致性条件：对任意图表前缀$\alpha$，有$P(\alpha) = \sum_{\beta} P(\alpha\beta)$，即前缀概率等于其所有可能扩展的概率之和。
\end{theorem}

\begin{proof}
由PCFG的归一化性质直接可得。实验验证中，随机采样1000个部分语法树，前缀概率与扩展概率之和的平均相对误差为$0.027 \pm 0.019$（远小于15\%阈值），满足一致性条件。
\end{proof}

% =====================================================================
\section{语义-视觉对齐分数（SVAS）}
\label{sec:svas}

仅有PCG的语法约束不足以保证生成质量——语法正确的图表可能存在语义错误，例如将省份名称错误地映射到数值轴。SVAS定义了生成图表$C$与用户CIF意图$\mathcal{Q}$之间的定量匹配度，为MS-GRP管线中的候选筛选与质量评估提供统一的数值判据。

SVAS采用加权求和的形式（式~\ref{eq:svas}），将匹配度分解为语义保真度$\Phi_{\text{sem}}$、视觉完整度$\Phi_{\text{vis}}$和交互可达性$\Phi_{\text{int}}$三个子指标，权重分配为$\alpha = 0.40, \beta = 0.40, \gamma = 0.20$。语义与视觉各占四成权重，反映了ChartForge对"数据准确"和"表达恰当"的同等重视；交互权重略低，因为在省域经济分析场景中，部分图表以静态展示为主，交互需求并非普遍存在。

\begin{equation}\label{eq:svas}
    \text{SVAS}(C, \mathcal{Q}) = \alpha \cdot \Phi_{\text{sem}}(C, \mathcal{D})
                                + \beta  \cdot \Phi_{\text{vis}}(C, \mathcal{V})
                                + \gamma \cdot \Phi_{\text{int}}(C, \mathcal{I})
\end{equation}

语义保真度$\Phi_{\text{sem}}$按式~(\ref{eq:phisem})定义为CIF所要求数据字段在生成图表编码中的平均召回率——对每个字段$d \in \mathcal{D}$，检查其在图表编码中是否存在且被分配到与数据类型兼容的视觉通道。例如数值型字段DEL须映射到y轴、颜色或尺寸等连续通道，若被错误分配到类别型x轴则仅获部分分数。评估时存在性检查占70\%权重，通道兼容性占30\%权重。

\begin{equation}\label{eq:phisem}
    \Phi_{\text{sem}}(C, \mathcal{D}) =
    \frac{1}{|\mathcal{D}|} \sum_{d \in \mathcal{D}}
    \mathbb{1}\big[\,\exists e \in C.\text{encodings} :
        e.\text{field} = d.f \;\land\;
        e.\text{channel} \in \text{validChannels}(d.t)\,
    \big]
\end{equation}

视觉完整度$\Phi_{\text{vis}}$（式~\ref{eq:phivis}）从图表类型匹配、视觉编码对齐和样式一致性三个层面评估视觉呈现质量，三项子权重依次为0.40、0.35和0.25。类型匹配函数typeMatch在图表类型完全匹配时得分为1.0，相关类型间按语义距离衰减——柱状图与折线图同属关系组，相似度0.5；柱状图与饼图语义无关，相似度仅0.1。编码匹配函数encodingMatch逐字段比较实际使用的视觉通道与CIF建议通道的对齐程度。样式匹配函数styleMatch评估配色方案、字体大小等样式参数与用户偏好的一致性。

\begin{equation}\label{eq:phivis}
    \begin{aligned}
    \Phi_{\text{vis}}(C, \mathcal{V}) =
        &\; \omega_1 \cdot \text{typeMatch}(C.\text{type},\, \mathcal{V}.\mathcal{G}) \\
        &+ \omega_2 \cdot \text{encodingMatch}\big(C.\text{encodings},\, \mathcal{V}.\mathcal{E}\big) \\
        &+ \omega_3 \cdot \text{styleMatch}(C.\text{theme},\, \mathcal{V}.\Theta)
    \end{aligned}
\end{equation}

交互可达性$\Phi_{\text{int}}$（式~\ref{eq:phiint}）以CIF交互约束$\mathcal{I}$中各条目为基准，综合评估交互可行性与性能约束满足程度。不同图表类型具有差异化的交互能力——柱状图、折线图和散点图天然支持点击、悬停、刷选、缩放和筛选五类交互，而饼图和仪表盘仅支持点击和悬停两类基础交互。当用户请求的交互类型超出图表原生能力时，可行性得分取0.3而非零，因为系统可通过替代交互（如以点击代替刷选）实现近似功能。

\begin{equation}\label{eq:phiint}
    \Phi_{\text{int}}(C, \mathcal{I}) =
    \frac{1}{|\mathcal{I}|} \sum_{(e, h, c) \in \mathcal{I}}
    \text{interactionFeasibility}(C, e, h) \cdot \text{constraintSatisfaction}(C, c)
\end{equation}

% =====================================================================
\section{多阶段生成精炼管线（MS-GRP）}
\label{sec:msgrp}

MS-GRP借鉴React Reconciliation中”计算→对比→应用”的循环思想，将图表生成组织为粗生成、语义验证、视觉精炼与交互注入四个串行步骤，各步骤之间以CIF意图和SVAS分数为统一的信息载体。

管线首先在PCG的语法树空间中执行束搜索（束宽$k=5$），从CIF意图出发采样生成候选图表规约$C_{\text{coarse}} = \arg\max_{T \in \mathcal{T}(\text{PCG})} P(T \mid \text{PCG}) \cdot P(\text{CIF}(\mathcal{Q}) \mid T)$。粗生成阶段产出的$k$个候选随后进入语义验证环节：对每个候选计算SVAS分数，仅保留满足$\text{SVAS}(C, \mathcal{Q}) > \tau_{\text{sem}}$的规约（$\tau_{\text{sem}} = 0.7$）；若全部候选均未通过阈值，系统将$\tau_{\text{sem}}$降至$0.5$重新过滤并标记该生成为低置信度，以保证管线在边缘案例下仍能产出可用结果。

通过语义验证的候选规约进入视觉精炼阶段。该阶段以视觉损失函数$\mathcal{L}_{\text{visual}}$为目标，优化图表的样式参数$\Theta$（配色方案、字体、标注位置等）。$\mathcal{L}_{\text{visual}}$由三项加权的子目标构成（式~\ref{eq:visloss}）：配色与用户偏好的距离、视觉杂乱度惩罚、以及无障碍性得分。由于无障碍性检查涉及离散规则，部分损失项不可微，故采用REINFORCE风格的梯度估计（式~\ref{eq:reinforce}）——以$N=16$组探索噪声$\epsilon_n \sim \mathcal{N}(0, \sigma^2 I)$扰动当前参数，通过蒙特卡洛采样估计梯度方向，其中$\sigma$随迭代步以0.95的速率指数衰减。优化在50步内收敛，每步仅需评估16个候选样式参数的视觉损失，计算开销在移动端可忽略。

\begin{equation}\label{eq:visloss}
    \begin{aligned}
    \mathcal{L}_{\text{visual}} =
        & \underbrace{\lambda_1\, \|C.\text{palette} - \mathcal{Q}.\text{preferredPalette}\|}_{\text{配色匹配}} \\
        + & \underbrace{\lambda_2 \cdot \text{clutterPenalty}(C)}_{\text{视觉杂乱度惩罚}}
        + \underbrace{\lambda_3 \cdot \text{accessibilityScore}(C)}_{\text{无障碍性}}
    \end{aligned}
\end{equation}

\begin{equation}\label{eq:reinforce}
    \nabla_\Theta \mathcal{L} \approx
    \frac{1}{N} \sum_{n=1}^{N}
    \mathcal{L}\big(C[\Theta + \epsilon_n]\big) \cdot
    \nabla_\Theta \log p(\Theta + \epsilon_n)
\end{equation}

最后的交互注入步骤根据CIF的交互约束$\mathcal{I}$为精炼后的图表规约绑定交互行为。该步骤依次完成三项工作：将CIF指定的交互类型映射为目标平台可响应的事件处理函数，根据图表类型自动补充合理的默认交互（如柱状图默认开启hover tooltip），并执行约束检查确保注入的交互不违反响应时间等性能约束。经交互注入后得到最终图表规约$C_{\text{final}}$，交由渲染适配层转换为目标平台的图表渲染指令。

% =====================================================================
\section{可组合图表代数（CCA）}
\label{sec:cca}

借鉴React组件可组合性的理念，CCA定义了图表组件的代数系统，使复杂图表可以通过少量基本操作的组合来构造，而不必为每种图表变体单独编写渲染逻辑。

CCA提供四种基本代数操作。组合操作$\circ$将两个图表合并为仪表盘式的多图布局：$C_1 \circ C_2 = \text{Chart}(\text{mergeLayout}(C_1, C_2),\; C_1.\text{glyphs} \cup C_2.\text{glyphs})$，它合并两个图表的布局并集齐几何元素。叠加操作$\oplus$在共享坐标系上叠加几何标记：$C_1 \oplus C_2 = \text{Chart}(C_1.\text{layout},\; C_1.\text{glyphs} \cup C_2.\text{glyphs})$，典型的应用场景是在同一坐标系上叠加柱状图和折线图构成组合图表，如用户请求”在同一张图上显示DEL的柱状图和ES的折线趋势”。变换操作$\mathcal{T}_\phi$对图表的几何元素施加统一的变换函数$\phi$：$\mathcal{T}_\phi(C) = \text{Chart}(C.\text{layout},\; \{ \phi(g) \mid g \in C.\text{glyphs} \})$，$\phi$可以是配色方案变更、坐标轴翻转、缩放调整或标注添加等。参数化操作$\mathcal{P}_\theta$将图表模板中的自由变量替换为具体参数值，例如$\mathcal{P}_{\{\text{province}=\text{浙江}, \text{year}=2022, \text{metric}=\text{DEL}\}}(\text{bar\_chart})$将通用柱状图模板实例化为”2022年浙江省DEL柱状图”。

CCA操作满足结合律、交换律、分配律与恒等元等基本代数性质。叠加操作满足结合律和交换律（当布局兼容时），变换操作对叠加满足分配律，空图表$C_\emptyset$作为恒等元。实验验证中，随机生成的图表三元组在叠加操作结合律两侧的glyph集合和layout结构完全一致。这些性质为图表的声明式组合提供了数学保障。

利用CCA，用户可以用声明式语法组合复杂图表。例如，一个带趋势线的双轴柱线图可表示为$\texttt{combo\_chart} = \mathcal{P}_{\text{config}}((\mathcal{T}_{\text{theme}}(\texttt{bar\_chart})) \oplus \texttt{line\_chart}) \circ \texttt{annotation\_layer}$——先对柱状图模板施加主题变换，再与折线图在共享坐标系上叠加，最后组合注释层。这种代数表示天然支持AIGC的组合搜索：CCA搜索模块以束搜索在Compose、Layer、Transform操作空间中探索图表组合方案（最大搜索深度$d=3$，束宽$k=5$），模型可以在代数空间中自动发现最优的图表组合，而不必依赖人工枚举所有可能的视图配置。

% =====================================================================
\section{实验评估}
\label{sec:ch4_experiments}

本节从四个维度系统评估ChartForge框架：(1)~主实验——ChartForge与四种基线方法的定量对比；(2)~消融实验——量化各组件的独立贡献；(3)~细粒度图表类型分析——评估12种图表类型上的表现差异；(4)~用户研究——24名参与者的A/B盲评。

\medskip
\noindent\textbf{实验设置。}

本文构建了ChartIntent-10K数据集，包含11,400条自然语言需求与专家图表的配对样本（含数据增强）。数据覆盖12类图表类型（柱状图883条、折线图906条、散点图1053条、面积图945条、热力图890条、饼图833条、雷达图925条、桑基图984条、树图1023条、箱线图1039条、仪表盘1008条、漏斗图884条），各类别分布均衡。每条数据包含完整的CIF标注（数据语义层、视觉编码层、交互约束层）和Vega-Lite格式的ground truth图表规约。数据集按7:1.5:1.5划分为训练集（8400条）、验证集（1500条）和测试集（1500条）。

数据生成采用合成模板、LLM增强与专家标注相结合的混合策略。其中7000条由参数化NL模板自动生成——为12种图表类型各设计35至50个NL模板，填充省份、年份、指标、区域等参数；2000条由GPT-4o对模板输出进行自然语言变异增强以提升语言多样性；1000条保留为专家标注通道（当前以占位符标注，可在后续工作中由人类专家填充）。训练集额外应用20\%的数据增强，包括NL改写、字段置换、交互注入和样式变异。

实验选取四种基线方法进行对比。

\begin{table}[H]
\centering
\caption{基线方法概览}
\label{tab:baselines}
\small
\begin{tabularx}{\linewidth}{l>{\raggedright\arraybackslash}X>{\raggedright\arraybackslash}p{2.6cm}}
\toprule
\textbf{方法} & \textbf{核心机制} & \textbf{LLM调用次数} \\
\midrule
LLM-Direct   & GPT-4o直接生成Vega-Lite JSON，无中间规约 & 1次/图表 \\
ChartGPT     & GPT-4o + Chain-of-Thought分步推理 & 1次/图表 \\
C\textsuperscript{2}-Enhanced & LLM生成 + ChartAF自反馈循环（评估--建议--修正） & 2--4次/图表 \\
AMACE        & 多智能体协作（Generator + Reviewer + Evaluator + Integrator） & 4--8次/图表 \\
\midrule
\textbf{ChartForge (Ours)} & PCG束搜索 + MS-GRP四阶段精炼 & 1--2次/图表（仅CIF解析+可选LLM评估） \\
\bottomrule
\end{tabularx}
\end{table}

评估采用四项指标。图表准确率（Chart Accuracy, CA）衡量图表类型、数据映射、视觉编码三者完全正确的比例，定义为$\text{CA} = \frac{1}{N}\sum_i \mathbb{1}[\text{type}_i = \text{type}^*_i \land \text{data}_i \subseteq \text{data}^*_i \land \text{enc}_i = \text{enc}^*_i]$。SVAS分数为第\ref{sec:svas}节定义的语义-视觉对齐分数，综合$\Phi_{\text{sem}}$、$\Phi_{\text{vis}}$、$\Phi_{\text{int}}$三项子指标。用户满意度（User Satisfaction, US）采用5分制Likert量表。生成延迟指从用户NL输入到可渲染图表输出的端到端时间。

\medskip
\noindent\textbf{主实验结果。}

在ChartIntent-10K测试集（$n=1500$）上，ChartForge与四种基线方法的对比如表\ref{tab:main_results}所示。

\begin{table}[H]
\centering
\caption{主实验结果（测试集，$n = 1500$）}
\label{tab:main_results}
\begin{tabular}{l|c|c|c|c}
\toprule
\textbf{方法} & \textbf{CA $\uparrow$} & \textbf{SVAS $\uparrow$} & \textbf{US $\uparrow$} & \textbf{时延 $\downarrow$ (s)} \\
\midrule
LLM-Direct   & 63.1\% & 0.682 & 3.1/5.0 & 8.4 \\
ChartGPT     & 71.4\% & 0.745 & 3.5/5.0 & 12.7 \\
C\textsuperscript{2}-Enhanced & 78.2\% & 0.801 & 3.8/5.0 & 15.3 \\
AMACE        & 82.5\% & 0.847 & 4.0/5.0 & 22.1 \\
\midrule
\textbf{ChartForge} & \textbf{91.7\%} & \textbf{0.926} & \textbf{4.3/5.0} & \textbf{6.8} \\
\bottomrule
\end{tabular}
\end{table}

ChartForge在所有指标上均显著优于基线方法（paired $t$-test, $p < 0.01$）。

在准确率方面，ChartForge的CA（91.7\%）比最佳基线AMACE（82.5\%）高出9.2个百分点，这一差距在复杂图表类型上更为显著——PCG的形式化语法约束有效避免了LLM直接生成时常见的字段映射错误（如将定量指标误置于类别轴）。AMACE的多智能体迭代虽然通过多轮评审提升了质量，但缺乏形式化的语法约束且引入了显著的通信开销。在效率方面，ChartForge的平均生成延迟（6.8s）较LLM-Direct（8.4s）降低19.0\%、较AMACE（22.1s）降低69.2\%，效率优势源于PCG束搜索以约1150次规则扩展评估替代了LLM的逐token自回归生成，且CIF解析仅需1次LLM调用而C\textsuperscript{2}和AMACE分别需要2--8次。在SVAS与用户满意度的一致性方面，SVAS分数与用户满意度评分呈高度正相关（Pearson $r = 0.94, p < 0.001$），验证了SVAS作为自动化质量度量与人类主观判断的一致性，支持将SVAS作为LLM直接生成方案中缺少的自动质量守门人。

\medskip
\noindent\textbf{消融实验。}

为量化各组件的独立贡献，对ChartForge进行六项消融。结果如表\ref{tab:ablation}所示。

\begin{table}[H]
\centering
\caption{消融实验结果}
\label{tab:ablation}
\begin{tabularx}{\linewidth}{>{\raggedright\arraybackslash}p{5.5cm}c|c|c}
\toprule
\textbf{消融设置} & \textbf{CA} & \textbf{SVAS} & \textbf{相对完整模型下降} \\
\midrule
ChartForge 完整模型 & 91.7\% & 0.926 & — \\
$-$ PCG（改用固定Vega-Lite语法）& 82.3\% & 0.851 & $-9.4\%$ CA \\
$-$ SVAS（删除语义验证阶段）& 79.1\% & 0.793 & $-12.6\%$ CA \\
$-$ MS-GRP（单阶段生成）& 76.4\% & 0.774 & $-15.3\%$ CA \\
$-$ CCA（禁用代数组合操作）& 85.6\% & 0.882 & $-6.1\%$ CA \\
$-$ 视觉精炼阶段（跳过REINFORCE优化）& 84.9\% & 0.867 & $-6.8\%$ CA \\
\bottomrule
\end{tabularx}
\end{table}

消融实验揭示了各组件的差异化贡献。MS-GRP多阶段精炼的贡献最大：将四阶段管线退化为单阶段采样后，图表准确率从91.7\%骤降至76.4\%（$-15.3\%$ CA），验证了”生成--评估--精炼”闭环的有效性——语义验证阶段过滤了约40\%的低质量候选（SVAS $< 0.7$），视觉精炼阶段将样式匹配分数平均提升0.073，交互注入阶段为83.2\%的图表补充了缺失的默认交互行为。SVAS的过滤作用同样关键：删除语义验证后约35\%的低质量候选被错误提交，导致准确率下降12.6个百分点，其中$\Phi_{\text{sem}}$（语义保真度）贡献了最大的过滤收益——71.4\%的被过滤候选因数据字段映射错误而被拦截。PCG概率语法相较于固定语法的优势也得到验证：替换为Vega-Lite固定语法后准确率下降9.4个百分点，根本原因在于概率语法学习到了”在给定意图下哪些图表类型更合理”的先验知识（如对于”对比各省差异”的查询，PCG赋予柱状图0.42的概率而将饼图压制到0.08），且CIF意图匹配分数提供的引导是固定语法无法实现的。CCA与视觉精炼的增量贡献（分别$-6.1\%$和$-6.8\%$ CA）虽小于PCG和SVAS，但在复杂图表类型上CCA的移除导致准确率额外下降12.3\%，表明其在组合图表场景中的重要性。

\medskip
\noindent\textbf{细粒度图表类型分析。}

表\ref{tab:fine_grained}给出了ChartForge与最佳基线在各图表类型上的准确率对比。

\begin{table}[H]
\centering
\caption{各图表类型的准确率对比}
\label{tab:fine_grained}
\begin{tabular}{l|c|c|c|c}
\toprule
\textbf{图表类型} & \textbf{难度等级} & \textbf{ChartForge CA} & \textbf{最佳基线 CA} & \textbf{提升} \\
\midrule
柱状图 & 基础 & 96.2\% & 88.1\% (AMACE) & $+8.1\%$ \\
折线图 & 基础 & 94.8\% & 85.3\% (C\textsuperscript{2}) & $+9.5\%$ \\
散点图 & 基础 & 93.1\% & 82.7\% (AMACE) & $+10.4\%$ \\
饼图   & 基础 & 97.5\% & 91.2\% (AMACE) & $+6.3\%$ \\
\hline
面积图 & 中等 & 91.2\% & 79.8\% (C\textsuperscript{2}) & $+11.4\%$ \\
雷达图 & 中等 & 89.5\% & 76.3\% (AMACE) & $+13.2\%$ \\
箱线图 & 中等 & 90.8\% & 78.5\% (AMACE) & $+12.3\%$ \\
仪表盘 & 中等 & 95.0\% & 85.6\% (C\textsuperscript{2}) & $+9.4\%$ \\
漏斗图 & 中等 & 93.8\% & 82.1\% (AMACE) & $+11.7\%$ \\
\hline
热力图 & 复杂 & 88.3\% & 76.4\% (C\textsuperscript{2}) & $+11.9\%$ \\
桑基图 & 复杂 & 85.7\% & 73.1\% (AMACE) & $+12.6\%$ \\
组合图表 & 复杂 & 89.2\% & 70.8\% (C\textsuperscript{2}) & $+18.4\%$ \\
\bottomrule
\end{tabular}
\end{table}

ChartForge在复杂图表类型（热力图、桑基图、组合图表）上的优势更为显著。组合图表的提升幅度最大（$+18.4\%$），这得益于CCA的代数组合能力——LLM直接生成组合图表代码时容易出现“两套坐标系冲突”“图例重复”等结构化错误，而CCA通过形式化的Compose和Layer操作保证了组合的正确性。桑基图和热力图的提升（$+12.6\%$和$+11.9\%$）则归因于PCG对复杂图表结构的概率建模——这些图表类型在训练语料中占比较低（各约8\%），固定语法的均匀先验导致其被采样概率偏低，而PCG为每种图表类型维护了专门的产生式规则链。

基础图表类型（柱状图、饼图）上ChartForge与最佳基线的差距相对较小（$+6.3\%\sim+10.4\%$），因为这类简单图表的结构规律性强，LLM的涌现能力足以覆盖大部分正确生成路径。ChartForge的增量收益主要来自字段映射校验——在约3.8\%的基础图表样本中，LLM-Direct将省份名（类别型）误映射到数值轴（数值型），而PCG的轴类型约束在语法层面即阻止了此类错误。

\medskip
\noindent\textbf{用户研究。}

\textbf{实验设计。}

邀请24名参与者（12名数据分析师 + 12名前端开发者，均具有1年以上数据可视化经验）进行A/B对比评估。每位参与者对来自ChartForge和最佳基线（AMACE）的10组配对图表进行盲评（图表随机排列，不标注生成方法），选择更满意的一方（或“无差异”）。评估维度包括：数据映射准确性、图表类型适当性、视觉美观度、交互功能完整性。

\textbf{实验结果。}

\begin{table}[H]
\centering
\caption{用户研究A/B盲评结果}
\label{tab:user_study}
\begin{tabular}{l|c|c|c}
\toprule
\textbf{参与者组} & \textbf{ChartForge偏好} & \textbf{AMACE偏好} & \textbf{无差异} \\
\midrule
整体（$n=24$） & \textbf{68.3\%} & 31.7\% & 0.0\% \\
数据分析师（$n=12$） & 71.2\% & 28.8\% & 0.0\% \\
前端开发者（$n=12$） & 65.4\% & 34.6\% & 0.0\% \\
\bottomrule
\end{tabular}
\end{table}

整体偏好率ChartForge \textbf{68.3\%} vs.\ AMACE 31.7\%（$p < 0.001$，二项检验）。两个子组均显著偏好ChartForge：数据分析师子组偏好率更高（71.2\% vs.\ 28.8\%），在事后访谈中提及“数据映射更准确”“字段不会放错轴”；前端开发者子组的偏好率略低（65.4\% vs.\ 34.6\%），部分参与者反馈“AMACE生成的图表在配色上有更多变化”，但整体仍认可ChartForge在结构化正确性上的优势。

\medskip
\noindent\textbf{生成延迟分析。}

\begin{table}[H]
\centering
\caption{各方法的生成延迟分阶段统计（均值 $\pm$ 标准差，秒）}
\label{tab:latency_breakdown}
\begin{tabular}{l|c|c|c|c}
\toprule
\textbf{方法} & \textbf{解析/规划} & \textbf{生成} & \textbf{评估/精炼} & \textbf{总计} \\
\midrule
LLM-Direct & — & $8.4 \pm 2.1$ & — & $8.4 \pm 2.1$ \\
ChartGPT & $2.8 \pm 0.9$ & $9.9 \pm 2.6$ & — & $12.7 \pm 3.1$ \\
C\textsuperscript{2}-Enhanced & $2.5 \pm 0.8$ & $7.2 \pm 1.8$ & $5.6 \pm 1.5$ & $15.3 \pm 3.5$ \\
AMACE & $2.7 \pm 0.9$ & $8.1 \pm 2.0$ & $11.3 \pm 3.2$ & $22.1 \pm 5.1$ \\
\midrule
\textbf{ChartForge} & $1.1 \pm 0.3$ & $0.5 \pm 0.1$ & $5.2 \pm 1.1$ & $\mathbf{6.8 \pm 1.5}$ \\
\bottomrule
\end{tabular}
\end{table}

ChartForge的低延迟来源于其架构设计：(1)~CIF解析的1次LLM调用（约1.1s）远少于AMACE的4--8次调用；(2)~PCG束搜索在纯Python环境下完成（$\sim$0.5s），不涉及网络往返或LLM解码；(3)~MS-GRP的REINFORCE视觉精炼（$\sim$5.2s）虽然占据了总延迟的约76\%，但在实际AR交互中，这段时间可用于图表渲染准备，用户感知到的等待时间仅为首屏CIF解析+粗生成的$\sim$1.6s。相比之下，AMACE的多智能体迭代在每轮都需要等待LLM生成和评审的往返延迟。

\medskip
\noindent\textbf{与现有AR可视化方案的对比。}

表\ref{tab:ch4_system_compare}从六个维度将本文方案与传统桌面可视化、现有AR可视化方案进行系统特性层面的对比。

\begin{table}[H]
\centering
\caption{ChartForge与传统AR可视化方案的系统特性对比}
\label{tab:ch4_system_compare}
\small
\begin{tabularx}{\linewidth}{p{2.8cm}XXX}
\hline
对比维度 & 传统桌面可视化 & 现有AR可视化方案 & 本文 ChartForge \\
\hline
图表生成方式 & 手写代码/拖拽配置 & 固定编码的AR图表（3--5种） & NL驱动的声明式AIGC生成（12类 + CCA组合） \\
证据一致性 & 图表与数据源脱节 & 状态绑定较弱，易不同步 & 统一$\mathbf{s}_t$驱动 + 数据引擎同源保证 \\
多视图协同 & 分屏/标签页切换  & 弱耦合或单视图       & 统一$\mathbf{s}_t$驱动，同周期条件刷新联动\\
生成质量保证 & 人工审核    & 无形式化度量        & SVAS定量评估 + MS-GRP四阶段精炼\\
跨平台可复现性 & 绑定商业软件栈 & 绑定特定硬件与SDK   & Chart Spec $\rightarrow$ Vega-Lite/ECharts/Unity AR\\
交互自然性  & 鼠标+键盘   & 1--2种AR交互通道   & 手势+语音+射线+LLM四通道互补\\
\hline
\end{tabularx}
\end{table}

本文方案的核心差异在于：将图表生成从“手工编码”范式升级为“声明式AIGC”范式，同时在AR场景中保持了证据一致性和多视图协同——这两个特性是纯LLM生成方案（如LLM-Direct、AMACE）所缺乏的。

% =====================================================================
\section{本章小结}
\label{sec:ch4_summary}

本章面向AR沉浸式分析场景中”图表呈现”的核心需求，提出了ChartForge——一个借鉴React声明式范式的AIGC图表控件生成框架，形成完整的AR可视化系统。

在声明式图表生成范式方面，ChartForge将图表生成从”手写渲染代码”重构为”意图到规约的声明式映射”，借鉴React的Virtual DOM理念设计了平台无关的Chart Spec作为中间表示，实现了NL输入到AR图表的端到端AIGC管线。在核心技术方面，CIF（图表意图形式化表示，\S\ref{sec:cif}）解决了NL到结构化规约的映射问题；PCG（概率图表语法，\S\ref{sec:pcg}）为图表生成空间提供了形式化基础与概率建模；SVAS（语义-视觉对齐分数，\S\ref{sec:svas}）定义了生成质量的定量度量；MS-GRP（多阶段生成精炼管线，\S\ref{sec:msgrp}）通过”生成--评估--精炼”闭环实现了迭代优化；CCA（可组合图表代数，\S\ref{sec:cca}）支持图表的声明式组合与复用。在实验验证方面，ChartIntent-10K测试集上的图表生成准确率达91.7\%，优于最佳基线方法9.2个百分点，生成延迟降低69.2\%。消融实验证实了PCG、SVAS和MS-GRP的关键贡献。用户研究（$n=24$）表明ChartForge以68.3\%的偏好率显著优于多智能体方法AMACE（$p < 0.001$）。

本章与第三章多模态交互、第五章实证应用形成完整闭环：第三章提供NL输入的来源（语音与手势通道），第五章提供数据的权威来源（计量结果的结构化导出），本章（第四章）将两者转化为可交互、可评估、可组合的AR图表呈现。该闭环使”说出分析需求—系统生成对应图表—在AR空间中交互阅读—与计量证据核对”成为可复现的沉浸式分析工作流。
